\section{Measuring Accuracy: Confusion, Sensitivity \& Specificity}
\label{sec:metrics}

Every detector in this report --- the shipped on-board freeze-index rule, the
glass-box RandomForest, and the FoGNet CNN --- ultimately hands down the same kind of
verdict: for each 4-second window it answers \emph{freeze} or \emph{no-freeze}. To
compare those detectors, and to decide where to set their trip lines, we need a
language for ``how good is this verdict?'' that survives contact with a hard clinical
fact: \textbf{freezes are rare}. A patient walks, stands and shuffles for minutes
between episodes that each last only seconds, so the no-freeze windows vastly
outnumber the freeze windows. This section builds that language from the bottom up ---
the confusion matrix, then the four rates derived from it --- and shows why the
single number most people reach for, raw accuracy, is the \emph{wrong} headline for
Cadence.

\begin{intuition}
Think of Cadence as a smoke alarm for the feet. Two ways to be right (stay quiet when
the kitchen is calm; sound when there is a fire) and two ways to be wrong (a false
alarm when you burn toast; silence during a real fire). A smoke alarm graded only on
``percent of seconds it was correct'' would score brilliantly by simply never going
off --- because almost every second of every day is fire-free. That is exactly the
trap a freeze monitor must not fall into.
\end{intuition}

\subsection{The confusion matrix: four ways a window can be scored}

Fix one detector and one trip line. Run it over a labelled recording and, window by
window, cross its \emph{prediction} against the \emph{ground truth}. Every window
lands in exactly one of four cells. We treat \emph{freeze} as the ``positive'' class,
because catching a freeze is the event we actually care about.

\begin{table}[H]
\centering
\caption{The $2\times2$ confusion matrix. Rows are what truly happened; columns are
what the detector said. The diagonal (TP, TN) is correct; the off-diagonal (FP, FN)
is the two distinct ways to be wrong --- and for a freeze monitor they cost very
different amounts.}
\label{tab:confusion}
\renewcommand{\arraystretch}{1.5}
\begin{tabular}{@{}c c >{\centering\arraybackslash}p{4.9cm} >{\centering\arraybackslash}p{4.9cm}@{}}
\toprule
& & \multicolumn{2}{c}{\textbf{Detector says}} \\
\cmidrule(l){3-4}
& & \textbf{freeze} (positive) & \textbf{no-freeze} (negative) \\
\midrule
\multirow{2}{*}{\rotatebox{90}{\textbf{\;Truth}}}
& \textbf{freeze}
& {\color{teal}\textbf{True Positive (TP)}} \par a real freeze, caught
& {\color{freeze}\textbf{False Negative (FN)}} \par a real freeze, \emph{missed} \\[4pt]
& \textbf{no-freeze}
& {\color{amber}\textbf{False Positive (FP)}} \par calm gait, false alarm
& {\color{teal}\textbf{True Negative (TN)}} \par calm gait, correctly ignored \\
\bottomrule
\end{tabular}
\end{table}

\begin{keyidea}
The two errors are \emph{not} symmetric in their consequences. A \textbf{false
negative} is a freeze the monitor failed to record --- and an unrecorded freeze is a
fall risk that the clinician never sees. A \textbf{false positive} is a moment of calm
gait flagged as a freeze --- annoying, and it muddies the log, but it hurts nobody.
Cadence is therefore tuned to push FN down even at the cost of a few more FP. This is
the single design decision behind the deployed freeze-index threshold of $1.815$
rather than the cleaner Youden optimum of $2.10$.
\end{keyidea}

\subsection{Four rates, and why accuracy is the wrong one}

From the four counts we build four rates. Each asks a different question of the
detector; each is just a ratio of cells in Table~\ref{tab:confusion}.

\begin{mathsbox}
\textbf{Sensitivity} (recall, true-positive rate) --- \emph{of the real freezes, what
fraction did we catch?}
\begin{equation}
\mathrm{sensitivity} = \frac{\mathrm{TP}}{\mathrm{TP}+\mathrm{FN}}
\label{eq:sens}
\end{equation}
\textbf{Specificity} (true-negative rate) --- \emph{of the calm-gait windows, what
fraction did we correctly leave alone?}
\begin{equation}
\mathrm{specificity} = \frac{\mathrm{TN}}{\mathrm{TN}+\mathrm{FP}}
\label{eq:spec}
\end{equation}
\textbf{Precision} (positive predictive value) --- \emph{when we shout ``freeze'', how
often are we right?}
\begin{equation}
\mathrm{precision} = \frac{\mathrm{TP}}{\mathrm{TP}+\mathrm{FP}}
\label{eq:prec}
\end{equation}
\textbf{Accuracy} --- \emph{of all windows, what fraction did we get right?}
\begin{equation}
\mathrm{accuracy} = \frac{\mathrm{TP}+\mathrm{TN}}{\mathrm{TP}+\mathrm{TN}+\mathrm{FP}+\mathrm{FN}}
\label{eq:acc}
\end{equation}
The crucial structural difference: sensitivity and specificity each look at \emph{one
row} of the matrix, so neither is diluted by how common the other class is. Accuracy
pools both rows --- and when one class dominates, it is dragged toward whichever rate
belongs to the majority.
\end{mathsbox}

That last sentence is the whole problem. Freezes occupy under a tenth of the windows
in a typical recording, so the no-freeze row dominates the accuracy sum. A detector
can score high accuracy while quietly failing at its one job.

\begin{figure}[H]
\centering
\includegraphics[width=0.74\linewidth]{accuracy_caveat}
\caption{The accuracy trap, made concrete. A degenerate ``do-nothing'' detector that
labels \emph{every} window no-freeze scores $0.90$ accuracy --- it beats our real
CNN's $0.83$ --- yet its sensitivity is exactly $0.00$: it catches not a single
freeze. Optimising for accuracy would actively reward the most useless possible
monitor. This is why Cadence reports sensitivity and specificity, never accuracy
alone.}
\label{fig:accuracy_caveat}
\end{figure}

\begin{takeaway}
For a rare event, \textbf{accuracy is a vanity metric}. A monitor that never flags a
freeze scores $\sim$\,$0.90$ accuracy and $0.00$ sensitivity --- useless, but it
``looks'' right. Cadence is judged on the pair (sensitivity, specificity): freezes
caught, and calm gait correctly left alone.
\end{takeaway}

\subsection{The real numbers: FoGNet under leave-one-subject-out}

The headline evaluation for the CNN upgrade uses \emph{leave-one-subject-out} (LOSO)
cross-validation on the public Daphnet FoG dataset: train on every subject but one,
test on the held-out subject, rotate through all subjects, then pool every held-out
window into one confusion matrix. LOSO is deliberately harsh --- the model is always
scored on a person it has never seen --- which is exactly the situation a new patient
strapping on the device represents.

\begin{figure}[H]
\centering
\includegraphics[width=0.6\linewidth]{confusion_matrix}
\caption{Pooled LOSO confusion matrix for FoGNet over $n=8982$ held-out windows.
Counts are TN $=6842$, FP $=1285$, FN $=248$, TP $=607$; the percentages are each
row normalised to its own total (so the top row reads as specificity, the bottom as
sensitivity).}
\label{fig:confusion}
\end{figure}

Plugging those four counts into
Equations~\eqref{eq:sens}--\eqref{eq:acc} gives the operating point of the shipped
CNN.

\begin{table}[H]
\centering
\caption{FoGNet pooled-LOSO metrics, computed from the counts in
Figure~\ref{fig:confusion}. Accuracy ($0.83$) sits \emph{below} the do-nothing
baseline of Figure~\ref{fig:accuracy_caveat} --- the clearest possible sign that
accuracy is not the metric to chase here.}
\label{tab:cnn_metrics}
\begin{tabular}{@{}l c l@{}}
\toprule
Metric & Value & From the counts \\
\midrule
Sensitivity & $0.71$ & $607 / (607+248)$ \\
Specificity & $0.84$ & $6842 / (6842+1285)$ \\
Accuracy    & $0.83$ & $(607+6842)/8982$ \\
\bottomrule
\end{tabular}
\end{table}

\begin{intuition}
Read the matrix as a story about one held-out patient cohort. Of every $100$ true
freeze windows, FoGNet logs about $71$ and lets $29$ slip past. Of every $100$
calm-gait windows it correctly stays quiet on $84$ and cries wolf on $16$. Neither
number is perfect --- this is a tiny model on a hard, noisy, real-world signal --- but
both are far above the chance line, and crucially neither is hidden by the dominance
of the calm-gait class.
\end{intuition}

\subsection{Sweeping the trip line: ROC, PR and Youden's J}

Every metric so far assumed a \emph{fixed} threshold. But a detector that emits a
score --- a freeze probability $\hat p$, or the raw freeze index --- has a whole
family of operating points: slide the threshold down and you catch more freezes
(sensitivity up) at the price of more false alarms (specificity down). Two curves
trace that trade-off across all thresholds at once. We illustrate them on the
threshold-free RandomForest glass-box, whose continuous score makes the sweep clean.

\begin{figure}[H]
\centering
\begin{minipage}[t]{0.49\linewidth}
\centering
\includegraphics[width=\linewidth]{roc_curve}
\end{minipage}\hfill
\begin{minipage}[t]{0.49\linewidth}
\centering
\includegraphics[width=\linewidth]{pr_curve}
\end{minipage}
\caption{Two threshold sweeps for the RandomForest glass-box. \textbf{Left:} the ROC
curve plots sensitivity against the false-alarm rate ($1-\text{specificity}$) as the
threshold slides; the dashed diagonal is a coin flip. Area under the curve,
$\text{AUC}=0.82$, summarises ranking quality in one number. \textbf{Right:} the
precision--recall curve, the honest view under $9.5\%$ freeze prevalence; its average
precision is only $\text{AP}=0.33$ against a no-skill floor of $0.095$. The gold star
marks the deployed operating point ($t\approx0.10$, sensitivity $0.75$, specificity
$0.76$) on both.}
\label{fig:roc_pr}
\end{figure}

\begin{keyidea}
ROC and PR answer different questions, and the gap between $\text{AUC}=0.82$ and
$\text{AP}=0.33$ is not a contradiction --- it is the rarity of freezes showing up
twice. \textbf{ROC} asks ``does the score \emph{rank} freezes above calm gait?'' and
is blind to how many of each there are, so it looks healthy. \textbf{Precision} asks
``when I shout freeze, am I right?'', and because true freezes are scarce, even a few
false positives drown the precision. PR is the metric that refuses to flatter a
rare-event detector.
\end{keyidea}

Why does the ranking look so much better than the precision? Because the score itself
\emph{does} separate the classes --- it just cannot do so perfectly. Figure~%
\ref{fig:score_hist} shows the two score distributions overlaid: calm-gait windows
pile up near zero, freeze windows are pushed to the right, but the tails overlap.
\emph{No} single threshold cleanly splits them; every choice trades one error for the
other.

\begin{figure}[H]
\centering
\includegraphics[width=0.78\linewidth]{score_hist}
\caption{Score separation for the glass-box. Each class is scaled to its own area so
the rare freeze class is visible. No-freeze windows (grey) cluster at low predicted
probability; freeze windows (red) spread right. The overlap is exactly the region
where any threshold must choose between a false positive and a false negative; the
gold line marks the deployed $t\approx0.10$.}
\label{fig:score_hist}
\end{figure}

\paragraph{Choosing the operating point.} Picking a threshold is picking one point on
the ROC curve. A principled, prevalence-free default is \textbf{Youden's $J$
statistic}, which rewards the threshold that buys the most sensitivity and specificity
together:
\begin{equation}
J = \mathrm{sensitivity} + \mathrm{specificity} - 1,
\label{eq:youden}
\end{equation}
maximised at the ROC point furthest above the chance diagonal. $J$ ranges from $0$
(useless) to $1$ (perfect), and it treats a missed freeze and a false alarm as equally
bad. For the on-board freeze-index rule on the worksheet capture, the $J$-optimal trip
line is $\FI > 2.10$ --- which on that capture gives perfect specificity ($1.00$) at
sensitivity $0.83$.

\begin{takeaway}
Cadence does \emph{not} ship the Youden optimum. Because a missed freeze can mean a
fall, the board biases toward sensitivity and trips at $\FI > 1.815$ --- buying back a
missed-freeze margin (sensitivity $0.83$) at the cost of one extra false alarm
(specificity $0.92$ instead of $1.00$). Youden tells you the balanced point; clinical
asymmetry tells you to lean past it. The headline is always the pair
\emph{(sensitivity, specificity)} --- never accuracy.
\end{takeaway}
